% theme: everyday_functions
\MajorQuestion{以下の問いに答えなさい。}
\SetRowSpace{3mm}

\Instruction{自動車の速さを $x\mathrm{km/h}$、ブレーキがきき始めてから停止するまでに進む距離を $y\mathrm{m}$ とする。$y$ は $x^2$ に比例し、速さが $30\mathrm{km/h}$ のときの制動距離は $6\mathrm{m}$ である。次の問いに答えなさい。}
\QQ{$y$ を $x$ の式で表しなさい。}{}
\QQ{速さが $60\mathrm{km/h}$ のときの制動距離を求めなさい。}{}
\QQ{制動距離が $54\mathrm{m}$ のときの速さを求めなさい。}{}
\QQ{速さを $30\mathrm{km/h}$ から $90\mathrm{km/h}$ にすると、制動距離は何倍になるか求めなさい。}{}
\QQ{速さを2倍にすると、制動距離は何倍になるか求めなさい。}{}
\FigProblemBlock{58mm}{fig/test_08_phone_charge}{\TextTall{ある電話サービスの1か月の料金は、通話時間が $15$ 分までは $1000$ 円である。$15$ 分をこえた後は、$5$ 分以下の端数も $5$ 分として数え、$5$ 分ごとに $100$ 円ずつ加算される。通話時間と料金の関係は右のグラフのようになる。次の問いに答えなさい。}}{
\QQ{通話時間が $15$ 分のときと $16$ 分のときの料金を、それぞれ求めなさい。}{}
\QQ{通話時間が $27$ 分のときの料金を求めなさい。}{}
\QQ{通話時間が $35$ 分のときと $36$ 分のときの料金を、それぞれ求めなさい。}{}
\QQ{料金が $1400$ 円となる通話時間 $x$ 分の範囲を、不等号を用いて表しなさい。}{}
\QQ{料金が $2500$ 円以下となる通話時間は、最長で何分か求めなさい。}{}
}
